\section{总结与延伸}

\begin{enumerate}
    \item RL 和 IL 各有优势，实际系统常结合使用
    \item Sim-to-real transfer 是 RL 用于真实机器人的核心流程
    \item Domain randomization 和 teacher-student 框架是弥合 sim-to-real gap 的关键技术
    \item RL 已在腿式运动和灵巧操作上取得突破性进展
    \item 人形机器人代表了下一代挑战
    \item 仿真器质量和奖励设计仍然是实际应用的瓶颈
\end{enumerate}

\begin{knowledgebox}{拓展阅读}
\begin{itemize}
    \item Kumar et al., ``Rapid Motor Adaptation (RMA),'' RSS 2021
    \item Tobin et al., ``Domain Randomization for Transferring Deep Neural Networks from Simulation to the Real World,'' IROS 2017
    \item OpenAI, ``Solving Rubik's Cube with a Robot Hand,'' 2019
    \item Rudin et al., ``Learning to Walk in Minutes Using Massively Parallel Deep RL,'' CoRL 2022
    \item Radosavovic et al., ``Real-World Humanoid Locomotion with RL,'' Science Robotics 2024
\end{itemize}
\end{knowledgebox}

\end{document}
